% Figure: the modulation transfer function (MTF) of a diffraction-limited
% incoherent imaging system against spatial frequency, normalised to
% the incoherent cutoff. The diffraction-limited incoherent MTF is the
% normalised autocorrelation of the circular pupil,
% MTF(s) = (2/pi)[acos(s) - s sqrt(1-s^2)] for 0<=s<=1, falling to zero at s=1.
% Overlaid: the coherent transfer function, a hard step that holds at unity to
% the coherent cutoff (half the incoherent cutoff) then drops to zero. The
% comparison shows that coherent imaging passes its band perfectly but cuts off
% at half the frequency.
\begin{figure}[htbp]
\centering
\begin{tikzpicture}
\begin{axis}[
  width=13cm, height=7cm,
  xlabel={spatial frequency $s$ (units of the incoherent cutoff $f_c$)},
  ylabel={transfer magnitude},
  xmin=0, xmax=1.05, ymin=0, ymax=1.05,
  axis lines=left,
  xtick={0,0.25,0.5,0.75,1.0},
  ytick={0,0.25,0.5,0.75,1.0},
  tick label style={font=\scriptsize},
  label style={font=\small},
  legend style={font=\scriptsize, at={(0.97,0.97)}, anchor=north east, draw=none},
  samples=200, domain=0:1,
]
% incoherent diffraction-limited MTF: (2/pi)(acos(s) - s sqrt(1-s^2))
% acos in pgf returns degrees; convert to radians with pi/180.
\addplot[engblue, very thick]
  {(2/pi)*( (pi/180)*acos(x) - x*sqrt(max(1-x^2,0)) )};
\addlegendentry{incoherent MTF (diffraction-limited)}
% coherent transfer function: unity to s=0.5 then zero. Draw as two segments.
\addplot[engred, very thick, domain=0:0.5] {1};
\addplot[engred, very thick, domain=0.5:1] {0};
\addlegendentry{coherent transfer (ideal)}
\draw[engred, dashed] (axis cs:0.5,0) -- (axis cs:0.5,1);
\node[engred, font=\scriptsize, anchor=south west] at (axis cs:0.5,0.55)
  {coherent cutoff $f_c/2$};
\node[engblue, font=\scriptsize, anchor=north east] at (axis cs:0.98,0.06)
  {incoherent cutoff $f_c$};
\end{axis}
\end{tikzpicture}
\caption[Coherent and incoherent transfer functions]{Spatial-frequency
response of a diffraction-limited imaging system with a circular pupil. The
incoherent modulation transfer function is the normalised autocorrelation of
the pupil, $\mathrm{MTF}(s) = (2/\pi)\,[\arccos s - s\sqrt{1-s^2}]$, falling
smoothly to zero at the incoherent cutoff $f_c = \mathrm{NA}/\lambda$ (in the
one-sided convention; $2\,\mathrm{NA}/\lambda$ for the full two-sided band).
The system passes coarse detail at high contrast and washes fine detail toward
the cutoff, the gentle roll-off every photographic and microscope system
shows. Coherent illumination instead passes its whole band at unit amplitude
and then stops dead at half the incoherent cutoff: a sharper transfer, but to
a lower frequency, which is why coherent imaging resolves edges with ringing
and incoherent imaging resolves finer periodic detail. The cutoff frequency,
set by the numerical aperture over the wavelength, is the spatial-frequency
form of the same diffraction limit the Airy disc states in the spatial
domain.}
\label{fig:vol04ch11:mtf-cutoff}
\end{figure}
